\chapter{Conclusioni e Sviluppi Futuri}
\section{Sintesi dei risultati}
Per questo elaborato è stato costruito un benchmark comparativo della letteratura sull'action anticipation con l'obiettivo di rendere i confronti tra modelli più equi e metodologicamente fondati rispetto a quelli attualmente esistenti. A partire da una raccolta di oltre 35 modelli proposti tra il 2015 e il 2024, sono stati individuati quattro gruppi di confronto in cui tutti i modelli condividono dataset, split, orizzonte di anticipazione, metrica e paradigma di task. L'analisi critica dei risultati ha permesso di identificare i fattori che contribuiscono maggiormente al progresso delle performance e le principali criticità metodologiche del campo. \\
Sul piano dei risultati quantitativi, i modelli migliori per ciascun gruppo sono: S-GEAR \cite{diko2024semantically} su EK-55 Val con 43.2\% Top-5 Action Accuracy, PlausiVL \cite{mittal2024can} su EK-100 Val con 27.6\% Mean Top-5 Action Recall, InAViT \cite{roy2024interaction} su EGTEA Gaze+ con 67.8\% Top-1 Action Accuracy, e ANTICIPATR \cite{nawhal2022rethinking} su Breakfast Actions con 37.4\% MoC all'orizzonte 20\%$\to10$\%. \\
Sul piano dell'analisi critica, emergono cinque osservazioni principali.  \\
La multimodalità contribuisce in modo consistente ma non uniforme: il tipo di \textit{cue} conta più della quantità, e la localizzazione esplicita delle interazioni mano-oggetto si rivela il contributo più determinante su dataset egocentrici come EGTEA Gaze+. L'architettura Transformer costituisce di per sé un vantaggio rispetto alle LSTM senza assunzioni strutturali specifiche per il task, come dimostrato dal caso di AVT \cite{girdhar2021anticipative}. La strategia di training può valere più dell'architettura, come dimostrato da DCR \cite{xu2022learning}, che porta un guadagno di circa 6 punti su EK-55 senza modificarne il modello di base. La finestra di osservazione ha un impatto inferiore alle aspettative nei gruppi STA, mentre l'orizzonte di anticipazione ha un effetto diretto e quantificabile nel gruppo LTA. Infine, i ranking tra dataset diversi mostrano bassa correlazione, confermando che EK-55, EK-100, EGTEA Gaze+ e Breakfast Actions misurano capacità distinte e non intercambiabili.
\section{Limiti del lavoro}
Il presente elaborato mostra alcuni limiti che è opportuno dichiarare esplicitamente. \\
Il primo riguarda la copertura della letteratura: la raccolta si è basata principalmente sui survey di Lai et al. \cite{lai2024human}, Hu et al.\cite{hu2022online} e Shuvo et al. \cite{shuvo2025deep}, dunque modelli rilevanti pubblicati dopo il 2024 o non coperti da questi lavori potrebbero non essere stati inclusi.\\
Il secondo riguarda i gruppi di confronto: i quattro gruppi individuati coprono solo una parte dei protocolli di valutazione presenti in letteratura. In particolare, non sono stati costruiti gruppi per THUMOS-14 e TVSeries, dataset rilevanti per il task STA, per insufficienza di modelli con task primario STA valutati su questi dataset con protocollo uniforme.

\section{Sviluppi futuri}
L'analisi condotta apre diverse direzioni di sviluppo che potrebbero estendere e approfondire il lavoro presentato.
La prima e più diretta è l'estensione del benchmark a nuovi modelli e dataset. Il campo dell'action anticipation è in rapida evoluzione, emergono con frequenza crescente modelli basati su Large Language Models e Vision-Language Models, e la loro inclusione in un benchmark sistematico richiederebbe un aggiornamento periodico dei gruppi di confronto. In particolare, il Gruppo 2 (EK-100) potrebbe essere arricchito con modelli più recenti che attualmente non riportano risultati sullo specifico split Val.
La seconda direzione riguarda la standardizzazione del protocollo di valutazione. I problemi di riproducibilità identificati in \ref{riprod} suggeriscono la necessità di un benchmark unificato con \textit{backbone} fissi, finestre di osservazione standardizzate e split dichiarati in modo uniforme. Un contributo futuro potrebbe consistere nell'effettuare un riesame di un sottoinsieme di modelli in condizioni completamente controllate, per isolare il contributo architetturale da quello del \textit{backbone} e della configurazione di input.





